\chapter{核心术语表}

\begin{longtable}{p{0.22\textwidth}p{0.68\textwidth}}
\toprule
术语 & 简释\\\midrule\endhead
\multicolumn{2}{l}{\textbf{知识、分析与方法}}\\
事实性 & 若主体知道命题，则命题为真；这不等于主体永不犯错。\\
命题知识 & 以“主体知道命题$p$”表达的知识，与亲知和能力之知相区分。\\
亲知 & 对人、地点或对象的熟悉性认识；其能否还原为命题知识有争议。\\
能力之知 & 知道怎样完成某项活动；能否仅由关于规则的命题知识构成有争议。\\
得到确证的真信念 & 由真、信念和确证三条件构成的知识分析，通常缩写为JTB。\\
必要条件 & 目标状态成立时必须满足的条件；反例要显示目标成立而条件不成立。\\
充分条件 & 条件满足便保证目标成立的条件；反例要显示条件满足而目标不成立。\\
反例 & 满足理论前提却不满足结论，或满足分析条件却不属于目标概念的案例。\\
最小对比 & 只改变一个认识上相关变量的成对案例，用来定位判断变化来源。\\
概念分析 & 研究概念的条件、结构和推论关系，可有描述、解释或改良目标。\\
概念工程 & 为解释或实践目标评估并改造现有概念，而非只描述日常用法。\\
知识优先 & 把知识视为基本认识状态，再用它解释证据、断言或行动的方案。\\
\multicolumn{2}{l}{\textbf{怀疑论与认识运气}}\\
怀疑论 & 对某类或几乎全部知识、确证或理性提出系统挑战的立场或方法。\\
局部怀疑论 & 只针对记忆、道德或某特定来源的怀疑，不直接推广到全部领域。\\
全局怀疑论 & 试图挑战几乎全部经验知识的怀疑方案，如完美模拟情景。\\
可错论 & 知识允许主体在相同类型根据下存在某些错误可能，不要求绝对确定。\\
不可错论 & 把排除相关错误可能或不可犯错性作为知识条件的强立场。\\
闭合原则 & 知识在主体掌握并完成的已知蕴涵推理下扩展的原则。\\
相关替代 & 在当前背景、证据或语境中必须排除的错误可能，而非所有逻辑可能。\\
摩尔式回应 & 以普通外部世界知识为前提，反向否定怀疑前提的论证策略。\\
语境主义 & “知道”归属句的适用标准随归属语境变化。\\
铰链命题 & 使怀疑、证据和求证实践得以运作的背景承诺或规则性命题。\\
判断悬置 & 在势均力敌或证据不足时暂不肯定也不否定，常以希腊词epoché指称。\\
欠决定 & 同一证据与多个竞争理论相容，证据本身不足以唯一选择其一。\\
标准问题 & 如何在“哪些信念是知识”与“什么标准能判定知识”之间避免循环的问题。\\
认识运气 & 信念以当前方式形成时很容易为假，真值与认识路径偶然相交。\\
介入运气 & 干扰因素进入信念形成链，并因另一因素抵消而碰巧成功。\\
环境运气 & 实际形成过程正常，但周围充满相似错误机会，如假谷仓环境。\\
可能世界 & 表示事情本可怎样的反事实模型；理论须说明哪些世界与实际世界相关或邻近。\\
敏感性 & 若命题为假，主体不会以相关方式相信它。\\
安全性 & 主体以相关方式相信时，在近邻情形中不容易为假。\\
因果偏离 & 事实虽进入信念的因果链，但其作用方式偶然或不适当，因而不足以构成知识。\\
认知德性 & 稳定导向真理的能力，或开放、认真、谦逊等认识人格品质。\\
成就 & 成功因为主体或系统的相关能力而取得，而非能力仅偶然参与。\\
\multicolumn{2}{l}{\textbf{确证、理由与结构}}\\
确证 & 使相信在认识意义上得到支持或合宜的地位。\\
命题确证 & 主体拥有支持命题的根据，不要求实际基于它相信。\\
信念确证 & 信念以适当方式基于主体拥有的根据。\\
基于关系 & 理由在信念产生或维持中实际发挥认识作用的因果、反事实或规范联系。\\
击败 & 新信息或事实撤销、削弱原有证据支持或信念认识地位的作用。\\
击败者 & 撤销或削弱原有支持的信息或事实。\\
反驳型击败者 & 提供非$p$证据、与原信念方向相反的击败信息。\\
削弱型击败者 & 不直接支持非$p$，却破坏原证据与$p$之间联系的信息。\\
误导性击败者 & 内容虽假，却因当时可信而合理地降低主体信心的信息。\\
高阶证据 & 不直接支持$p$或非$p$，而评价主体能力或证据处理质量的证据。\\
基础主义 & 某些非推论信念提供无需其他信念传递的初始确证。\\
融贯主义 & 确证来自信念系统的一致、解释联系和相互支持。\\
无限主义 & 认为理由链可以非循环地无限延伸，不以基础信念或循环终止。\\
内在主义 & 确证的决定因素必须以适当方式对主体反思可及。\\
外在主义 & 确证或知识可部分取决于主体无法反思把握的可靠关系。\\
可靠主义 & 信念的认识地位取决于产生或维持它的可靠过程。\\
证据主义 & 主体的信念态度应与其证据相称，争点包括什么构成主体证据。\\
一般性问题 & 同一信念形成事件属于多个过程类型，可靠主义需说明评价哪一类型。\\
实践侵入 & 实践利害或错误成本影响知识或确证门槛。\\
\multicolumn{2}{l}{\textbf{认识来源}}\\
知觉 & 通过感官和相关处理与对象、性质或事件发生认识联系的来源。\\
直接实在论 & 正常知觉直接呈现世界对象和性质，不以内部对象作为首要中介。\\
间接实在论 & 主体直接把握经验或感觉材料，再由其推知外部对象的立场。\\
意向论 & 把知觉经验理解为以表征内容呈现世界、具有准确条件的状态。\\
析取论 & 真知觉与不可区分幻觉不属于决定认识地位的同一根本种类。\\
认知渗透 & 信念、愿望或期望影响知觉经验内容，而非只影响后续判断。\\
知觉学习 & 训练与反馈使主体获得新的辨别、注意或分类能力。\\
记忆保存论 & 记忆主要保存先前获得的知识或确证，本身不通常创造新根据。\\
记忆生成论 & 当前记忆状态可整合痕迹并产生新的认识支持，而非仅机械保存。\\
来源监控 & 区分某内容来自亲历、证言、想象或后来阅读的记忆能力。\\
内省 & 主体获得自己当前体验、态度或心理状态认识的方式。\\
先验确证 & 最终确证不依赖特定感官经验；仍可需要经验来学习概念和语言。\\
归纳 & 从已观察样本或过去规律支持未观察实例和未来情况的可错推理。\\
证言 & 通过沟通把某内容呈现为真，并使他人可能据此形成信念。\\
证言还原论 & 证言权利最终需由知觉、记忆、归纳等非证言来源支持。\\
证言非还原论 & 理解看似真诚断言可给予可被击败的初步接受权利。\\
认识同侪 & 在相关证据、能力和认真程度上大致相当的分歧者。\\
专家元证据 & 非专家用来评估专家的领域训练、记录、透明度、利益和纠错表现。\\
来源谱系 & 信息从原始接触、记录、编辑、转述到听者的生成和传播结构。\\
\multicolumn{2}{l}{\textbf{形式认识论与决策}}\\
置信度 & 主体对命题为真的程度性信心，常用概率表示。\\
条件概率 & 在条件$B$成立时$A$的概率，记作$P(A\mid B)$；方向不可任意交换。\\
基础率 & 在获得个案证据前，目标状态在相关人群或参照类中的比例。\\
似然 & 假设成立时观察到某证据的概率，用来衡量证据区分假设的能力。\\
贝叶斯因子 & 两个假设对同一证据似然之比；与先验赔率结合得到后验赔率。\\
校准 & 某置信档位的长期发生频率与报告概率相匹配的性质。\\
分辨力 & 预测系统能给实际发生与未发生事件不同概率，而非总报基础率。\\
条件独立 & 给定假设后，两项证据的出现概率不再相互影响。\\
恰当评分规则 & 使主体按真实置信报告能最小化期望损失的概率评分方式。\\
期望效用 & 将行动在各可能状态的效用按状态概率加权后的总评价。\\
信息价值 & 新信息因可能改变更好行动而带来的期望收益，需扣除获取成本。\\
\multicolumn{2}{l}{\textbf{认识价值、社会与数字环境}}\\
理解 & 对解释、因果或依赖关系的把握，而非孤立事实的累积。\\
沼泽化问题 & 若真理是唯一最终价值，真信念到手后可靠过程的价值似乎被淹没。\\
智慧 & 识别什么值得知道、整合理由、承认边界并在不确定中适当行动的能力。\\
认识能动性 & 主体提出问题、选择证据、更新信念和塑造认识环境的能力。\\
社会认识论 & 研究个人的社会依赖以及群体、制度和网络认识表现的领域。\\
群体知识 & 由聚合、共同接受或功能结构构成、不能简单等同成员知识总和的知识。\\
联合承诺 & 成员通过程序共同接受并据此行动的一种群体态度结构。\\
认识不正义 & 主体作为知识提供者或理解者的地位因权力和偏见受损。\\
证言不正义 & 身份偏见使说话者获得不当的可信度亏欠。\\
诠释不正义 & 共同解释资源的结构性缺口使某些经验难以理解和表达。\\
制度客观性 & 通过公开标准、异质批评、复现和可修订程序降低个体偏差。\\
信息气泡 & 网络结构使主体遗漏相关来源，不必包含对外部来源的主动贬低。\\
回音室 & 内部网络同时削弱外部来源可信度，使反证被重新解释为敌对证明。\\
认识自主 & 识别根据、分配信任、响应反证并管理认识依赖的能力。\\
关系性自主 & 把自主理解为在支持或支配关系中形成并管理依赖的能力与条件。\\
自动化偏差 & 人过度接受自动系统输出、忽略相反信息或放弃独立判断的倾向。\\
算法厌恶 & 因一次显眼机器错误而拒绝总体表现较好的工具，未按任务校准依赖。\\
虚构 & 生成系统自信产生无来源支持或错误内容的风险；不预设故意欺骗。\\
检索增强生成 & 将外部文档提供给生成模型的流程，可改善来源但不保证综合正确。\\
内容溯源 & 记录文件创建、编辑、签名和传播历史的机制；不等于内容事实为真。\\
人机监督 & 人对系统输出进行审查、覆写和升级的安排，须有信息、时间与权限。\\
\multicolumn{2}{l}{\textbf{比较认识论}}\\
内部解释 & 在比较前依文本语境说明概念、论证与实践功能，不预设现代对应物。\\
结构比较 & 选择明确问题轴比较推论关系、错误机制或规范角色，而非只比较词形。\\
规范建构 & 研究者借多种传统资源提出新方案，并明确为自己的理论责任。\\
平行关系 & 不主张历史传播，只说明不同传统独立出现可比较的问题或结构。\\
历史同源 & 由传播、共同文本或接触证据支持的思想关联，不由年代先后单独证明。\\
概念簇 & 用对象、推论、反例和实践后果描述一组相关用法，避免一对一翻译。\\
厚概念 & 同时包含描述与评价、能力或实践要求的概念，不能轻易还原为薄命题状态。\\
不可通约 & 两种框架缺少无损共同尺度；它可揭示边界，不等于完全无法比较。\\
双向批评 & 比较双方都可改变问题并接受对方压力，而非一方提供唯一裁判标准。\\
\bottomrule
\end{longtable}
